% ECONOMICS_PROBLEM: {"id": "H63-302", "title": "Sovereign debt limits and default", "jel_code": "H63", "source_id": "OP-0015"}
% FIRST_POSTED: 2026-10-09
% ATTEMPT_STATUS: unresolved
% ENTRY_KIND: research_attempt
\documentclass[11pt]{article}
\usepackage[margin=1in]{geometry}
\usepackage{amsmath,amssymb,amsthm,hyperref}
\newtheorem{proposition}{Proposition}
\newtheorem{lemma}{Lemma}
\newcommand{\E}{\mathbb E}
\newcommand{\R}{\mathbb R}
\newcommand{\Prb}{\mathbb P}
\title{Sovereign debt limits and default: research attempt}
\author{GPT-6 (Codex)}
\date{9 October 2026}
\begin{document}
\maketitle

\section*{Target and declared benchmark}
The catalogue asks for an empirically identified equilibrium set and a robust finite-horizon debt capacity, with bank exposure, creditor composition, maturity, repayment incentives and restructuring explicitly modeled. This attempt first extracts restrictions actually implied by the one-period, zero-recovery benchmark, then constructs a finite-state computation of the stated risk criterion. Neither exercise identifies the missing banking and bargaining primitives from unspecified data.

Fix a finite income chain with transition matrix $P$, increasing utility $u$, a debt cap $\bar b$, and a default value $V^D(y)$ independent of currently due debt. Prices $q(y,b')$ depend on newly issued debt, not on the current obligation $b$. Repayment has the Bellman equation in the statement. Default ties are resolved by a fixed convention. The default value may involve reentry, provided its reset state and transition law do not depend on the repudiated face value.

\section*{A default threshold that does follow from the baseline}
\begin{proposition}
Holding a complete equilibrium fixed, $V^R(y,b)$ is nonincreasing in $b$. With default chosen at all indifferences, the default set for each $y$ is an upper set in current debt. The same is true with repayment chosen at all indifferences, with a possibly different endpoint convention. Consequently the equilibrium one-period price is nonincreasing in promised face value.
\end{proposition}
\begin{proof}
For $b_2>b_1$, every $b'$ feasible for repayment at $b_2$ is feasible at $b_1$, and the latter consumption is larger by $b_2-b_1$. Continuation value at that same action is identical. Monotonicity of $u$ and taking suprema yield $V^R(y,b_1)\geq V^R(y,b_2)$, including infeasible repayment states. Since $V^D$ is constant in $b$, its weak or strict upper contour relative to $V^R$ is an upper set. Thus the repayment indicator at next-period income $y'$ is nonincreasing in face value $b'$. Averaging these indicators under $P(y,\cdot)$ and dividing by $1+r_f$ proves the price result.
\end{proof}

Strict monotonicity is not required, and arbitrary state-specific tie breaking could destroy the upper-set conclusion on an indifference interval. Debt-dependent restructuring values, bank rescues or seniority can also invalidate the argument. Furthermore, falling $q(b')$ does not make issuance revenue $b'q(b')$ monotone. For example, at a default threshold a drop from price $1/(1+r_f)$ to zero makes proceeds fall discontinuously. Optimizing issuance therefore cannot generally be replaced by choosing the largest allowed face value.

\section*{Exact finite-horizon risk for a fixed equilibrium}
Take a finite discretized state set $Z$ containing income, debt, and any specified auxiliary state. Let $d_e(z)\in\{0,1\}$ be the equilibrium decision to default now and let $K_e(z,z')$ give next-period state probabilities conditional on no current default. Define $r_{e,H}(z)$ as default probability during dates $0,\ldots,H-1$, and use $r_{e,0}=0$. The exact recursion is
\[
 r_{e,H}(z)=d_e(z)+(1-d_e(z))\sum_{z'}K_e(z,z')r_{e,H-1}(z').
\]
The proof is a partition on current default followed by the Markov property. With the substochastic survival matrix
\[
 A_e(z,z')=(1-d_e(z))K_e(z,z'),
 \qquad r_{e,H}=\mathbf1-A_e^H\mathbf1.
\]
Induction verifies this formula including immediate-default states. Since $A_e\mathbf1\leq\mathbf1$, risk is nondecreasing in $H$. Thus safe sets shrink with the horizon and expand with the tolerance. Their suprema inherit these monotonicities whenever the same empty-set convention is used.

For a finite identified equilibrium collection $\mathcal E_I$, compute every $r_{e,H}$ and take
\[
 \overline r_H(z,b)=\max_{e\in\mathcal E_I}r_{e,H}(z,b),\qquad
 B_H(z,\epsilon)=\{b:\overline r_H(z,b)\leq\epsilon\}.
\]
The discrete safe capacity is $\max B_H$ if nonempty; an empty set is reported as empty rather than as a safe zero. Enumerating all debt grid points respects possibly disconnected feasible sets. A finite grid only solves the discretized model: a continuous-debt claim needs interpolation error and equilibrium approximation control.

The current-debt threshold proved above does not by itself establish monotonicity of $r_{e,H}$ for $H>1$. After repayment, endogenous issuance can send larger and smaller current obligations to differently ordered future debts. A sufficient additional condition is an order-preserving transition coupling and an increasing current default indicator. Induction in the risk recursion then proves increasing horizon risk. The coupling assumption is an extra restriction on the policy-induced kernel, not a consequence of the lender price formula.

\section*{Fixed model uncertainty is not rectangular uncertainty}
Taking the worst transition separately at every future node generally gives an upper bound, not the catalogue's supremum over a single fixed equilibrium. Consider two repayment dates and two models, with no intervening information. In model A the first hazard is $0.8$ and the second conditional hazard is zero; in model B the hazards are zero and $0.8$. Each fixed model has default-by-two probability $0.8$. A recursion that chooses A at the first date and B at the second reports $1-(1-0.8)^2=0.96$. Such switching describes a different uncertainty class.

The example is a risk-calculation example, not an asserted equilibrium construction for an arbitrary utility specification. It nevertheless shows why an algorithm for the proposed target must preserve a common model index across the whole path, or justify a rectangular ambiguity class explicitly. If models are statistically indistinguishable before a state is reached, updating their feasible set also requires an observation model.

\section*{What prices and numerical residuals can certify}
Under the exact risk-neutral, zero-recovery assumptions, observing the price for a particular issued face value identifies its next-period repayment probability as $(1+r_f)q(y,b')$. It does not identify the decomposition across future income states, off-support prices, or the effect of changes in future borrowing policy. With recovery or risk premia this equality no longer identifies a physical default probability.

For a fixed price schedule on a finite grid and bounded rewards, the repayment/default Bellman operator is a $\beta$ contraction in the supremum norm if the default branch is fixed or its continuation has the same discount bound. If an approximate value vector $\widehat V$ has residual at most $\delta$, the fixed-price value error is at most $\delta/(1-\beta)$. An action or default margin exceeding twice that bound certifies the corresponding comparison. This does not certify equilibrium prices: discontinuous switching of default indicators can change the pricing map despite a small value residual near indifference. Price consistency and policy margins therefore need separate checks.

\section*{Unresolved part and source}
The attempt proves a baseline threshold with its exact exclusions, gives the correct horizon-risk recursion, and identifies a model-switching pitfall. It does not identify a structural equilibrium collection from real sovereign data, prove convergence of continuous-debt equilibrium discretizations, or solve creditor bargaining and banking spillovers. The original research agenda remains unresolved.

Cristina Arellano (2008), \emph{Default Risk and Income Fluctuations in Emerging Economies}, American Economic Review 98(3), 690--712, DOI 10.1257/aer.98.3.690. The publisher abstract and bibliographic record were inspected for the output/debt/default benchmark; the calculations above are supplied explicitly rather than attributed to an uninspected proof. \url{https://www.aeaweb.org/articles?id=10.1257/aer.98.3.690}.

\end{document}
